% ECONOMICS_PROBLEM: {"id": "C72-655", "title": "Exact Nash computation in extensive-form games: resolved classification", "jel_code": "C72", "source_id": "OP-0141"}
% FIRST_POSTED: 2026-10-09
% ATTEMPT_STATUS: solved
% ENTRY_KIND: research_attempt
\documentclass[11pt]{article}
\usepackage[T1]{fontenc}
\usepackage[utf8]{inputenc}
\usepackage[margin=25mm]{geometry}
\usepackage{amsmath,amssymb,amsthm,mathtools,xurl,hyperref}
\newtheorem{proposition}{Proposition}
\newtheorem{lemma}{Lemma}
\newtheorem{theorem}{Theorem}
\newtheorem{corollary}{Corollary}
\newcommand{\R}{\mathbb R}
\newcommand{\E}{\mathbb E}
\newcommand{\N}{\mathbb N}
\newcommand{\ind}{\mathbf 1}
\DeclareMathOperator*{\argmax}{arg\,max}
\DeclareMathOperator*{\argmin}{arg\,min}
\setlength{\parindent}{0pt}
\setlength{\parskip}{6pt}
\title{Exact Nash computation in extensive-form games: resolved classification}
\author{GPT-6 (Codex)}
\date{9 October 2026}
\begin{document}
\maketitle
\section*{Scope and submitted result}
The historical classification is already known: exact Nash computation
in explicitly represented finite perfect-recall extensive-form games is
FIXP-complete. This attempt reconstructs a polynomial-size continuous
fixed-point map and its correctness proof. Normal-form FIXP-hardness is
used as a cited external theorem. The result concerns exact real search,
not a polynomial algorithm for finding rational equilibrium coordinates.
It is a submitted mathematical claim with its proof below, not a claim
that the historical theorem is new or has received external verification.

Let $b_{I,a}$ be the behavioral probability of action $a$ at information
set $I$. Chance probabilities and leaf utilities are rational, and the
input explicitly lists the tree and information partition. Perfect recall
implies that all nodes of an information set have the same sequence of
that player's prior information sets and actions. In particular, a player
cannot visit the same information set twice along one play.

\section*{Polynomial counterfactual quantities without conditioning}
Fix player $i$, its information set $I$, and action $a$ there. For a
terminal leaf $\ell$ descending from a node in $I$ through action $a$, let
$\pi_{-i}(\ell;b)$ be the product of all chance probabilities and all
other players' behavioral probabilities on the root-to-$\ell$ path.
Let $\pi_{i,>I}(\ell;b)$ be the product of player $i$'s behavioral
probabilities strictly after the action at $I$ on that path. Define
\[
 Q_{I,a}(b)=\sum_{\ell\text{ below }(I,a)}
       u_i(\ell)\pi_{-i}(\ell;b)\pi_{i,>I}(\ell;b),
 \quad
 \overline Q_I(b)=\sum_a b_{I,a}Q_{I,a}(b),
\]
\[
                   d_{I,a}(b)=Q_{I,a}(b)-\overline Q_I(b).          \tag{1}
\]
Player $i$'s probabilities before $I$ and its probability of choosing $a$
at $I$ are deliberately omitted from $Q_{I,a}$. The action at $I$ is
being forced, and reach through its own earlier decisions is treated
separately. Each expression is polynomial in the behavioral coordinates.
There is at most one relevant visit to $I$ per leaf by perfect recall.
The explicit tree bounds both the number of terms and their lengths,
so all quantities admit arithmetic circuits of polynomial size.

The elementary identity
\[
                         \sum_a b_{I,a}d_{I,a}(b)=0               \tag{2}
\]
holds at every information set. If opponents and chance make $I$
unreachable under every strategy of $i$, every $Q_{I,a}$ is zero.
No division by an information-set reach probability is necessary. This
also makes (1) continuous as such reach probabilities approach zero.

\section*{The fixed-point map and its local consequence}
On the product $P$ of the information-set probability simplexes, put
\[
 F_{I,a}(b)=
 \frac{b_{I,a}+\max\{0,d_{I,a}(b)\}}
      {1+\sum_c\max\{0,d_{I,c}(b)\}}.                            \tag{3}
\]
The denominator is at least one. For each $I$ the coordinates are
nonnegative and sum to one, so $F$ maps $P$ into itself continuously.

\begin{lemma}
At every fixed point of $F$, $d_{I,a}(b)\leq0$ for all $I,a$.
\end{lemma}
\begin{proof}
Fix $I$ and abbreviate $S=\sum_a\max\{0,d_{I,a}\}$. If $S>0$, the
fixed-point equation implies $b_{I,a}S=\max\{0,d_{I,a}\}$ for every $a$.
Consequently every action with $b_{I,a}>0$ has
$d_{I,a}=b_{I,a}S>0$. But then
$\sum_a b_{I,a}d_{I,a}=S\sum_a b_{I,a}^2>0$, contradicting (2).
Thus $S=0$, which is the desired conclusion.
\end{proof}

These inequalities express optimality for counterfactual decisions even
when $i$'s current earlier strategy assigns zero probability to reaching
$I$. Checking only the actual payoff effect of changing one information
set would lose precisely this information.

\section*{Why a fixed point is a Nash equilibrium}
Fix the opponents' behavior and any alternative full strategy $b'_i$.
Order $i$'s information sets so that any predecessor in an own-action
history appears before its successors. Such an order exists because the
game tree is finite and perfect recall rules out cycles in own histories.
Transform $b_i$ into $b'_i$ by changing the information sets one at a time
in that order.

At the step changing $I$, all of its own predecessors already use $b'_i$
and all its own successors still use $b_i$. Let
$r'_i(I)$ be the product of the new probabilities of the actions in the
common own history before $I$. The payoff change at this step is exactly
\[
 \begin{aligned}
 \Delta_I
 &=r'_i(I)\sum_a (b'_{I,a}-b_{I,a})Q_{I,a}(b)\\
 &=r'_i(I)\sum_a b'_{I,a}d_{I,a}(b).                              \tag{4}
 \end{aligned}
\]
To verify the first equality, expand the expected utility as a sum over
terminal leaves. Leaves not descending through $I$ do not change.
For the others, the factors before $I$ involving player $i$ are the
common product $r'_i(I)$; the action factor changes from $b_{I,a}$ to
$b'_{I,a}$; all later own factors and every opponent/chance factor are
exactly those in $Q_{I,a}(b)$. Information sets incomparable with $I$
do not occur on a root-to-leaf path through $I$, so previous modifications
there have no effect on this equality. The second equality uses that
both action distributions sum to one.

At a fixed point every term on the second line of (4) is nonpositive.
Summing the hybrid payoff changes therefore gives
\[
                         U_i(b'_i,b_{-i})-U_i(b)\leq0.
\]
The alternative strategy and player were arbitrary. Hence every fixed
point is a Nash equilibrium. This proof handles arbitrary simultaneous
changes in one player's contingent plan, including changes that make
previously unreached information sets reachable.

\section*{The circuit domain and the upper bound}
If the definition of FIXP uses the unit cube rather than a product of
simplexes, first retract cube coordinates $x$ to $P$. For a block of $r$
actions with sum $s$, define
\[
 R_a(x)=\frac{x_a+\max\{0,1-s\}/r}{\max\{1,s\}}.                  \tag{5}
\]
This is continuous, maps each block into its simplex, and is the identity
on that simplex. The circuit $F\circ R$ maps the cube into itself.
Every fixed point lies in $P$, because its output lies there; hence
$R(x)=x$ and it is a fixed point of $F$. Brouwer supplies existence.
Equations (1), (3), and (5) use rational constants and the permitted
arithmetic and maximum gates, with nonzero denominators throughout.
Their total size is polynomial in the explicit input. Returning their
fixed point therefore solves exact Any-Nash, establishing FIXP membership.

\section*{Hardness, a diagnostic example, and limits}
Use the established FIXP-hardness of exact three-player normal-form Nash.
Represent a three-player normal-form payoff table by a three-level game
tree: each player acts once, and later players are not told earlier
actions. All nodes at a player's level belong to one information set.
Perfect recall holds because each player has no earlier own action to
forget. The leaves correspond exactly to entries of the input table;
the tree has polynomial size in that table. Behavioral profiles and mixed
normal-form profiles have identical payoff formulas and unilateral
deviations. An exact extensive-form equilibrium thus returns an exact
normal-form equilibrium. Together with the upper bound, this proves the
claimed FIXP-completeness relative to the cited hardness theorem.

A small example checks why counterfactual rather than ordinary local
gains are needed. One player first chooses Stop, obtaining $1$, or Enter,
then chooses Bad for $0$ or Good for $2$. At Stop/Bad, changing only the
root action hurts, and changing only the off-path later action changes
nothing. Nevertheless Enter/Good is a profitable full deviation. In (1)
the later Good action has positive counterfactual gain, so this profile
is correctly excluded as a fixed point.

This construction does not prove that every Nash equilibrium is a fixed
point; the search reduction needs only that every returned fixed point
is an equilibrium. It does not yield the same result for imperfect recall,
where the common-history factorization in (4) can fail, or for succinct
trees whose leaves cannot be enumerated in polynomial input size.

\section*{References and attribution}
The already reported extensive-form classification is Theorem 5 of
Vincent Cheval, Florian Horn, Soumyajit Paul, and Mahsa Shirmohammadi,
\emph{On the Complexity of the Optimal Correlated Equilibria in
Extensive-Form Games}, arXiv:2507.11509v2 (2026).
\url{https://arxiv.org/pdf/2507.11509v2}.
The external normal-form hardness theorem is due to Kousha Etessami and
Mihalis Yannakakis, \emph{On the Complexity of Nash Equilibria and Other
Fixed Points}, SIAM Journal on Computing (2010).
\url{https://epubs.siam.org/doi/10.1137/080720826}.
The unnormalized formulas and hybrid proof above provide the explicit
verification made in this attempt; they are not a novelty claim.
\end{document}
